%%%PAGE 101 rot=0 legibility=good summary=光学知识提纲(思维导图式)：折射定律与全反射、色散、折射率理解及应用、综合应用公式、干涉衍射偏振、电磁波与相对论、考点
%%%BLOCK id=101-01 ch=15 sec=折射定律·全反射·色散 kind=summary alt=none cont=none y=0.03-0.22
①
% 原稿：左侧一个大括号并列下面三项
\begin{itemize}
\item 折射定律\quad $\dfrac{\sin\theta_1}{\sin\theta_2}=n$。
\item 全反射、光导纤维。
\item 光的色散\quad 材料对不同色光折射率不同\ $n_{\text{红}}$最小；$n_{\text{紫}}$最大，偏折程度最大
\end{itemize}
%%%END
%%%BLOCK id=101-02 ch=15 sec=折射定律和折射率的理解及应用 kind=summary alt=none cont=none y=0.24-0.31
②\ 光的折射。全反射、折射定律和折射率的理解及应用

$n$与介质、光的频率有关。同一色光在不同介质，$v$、$\lambda$ 不同但 $f$ 相同、光路可逆
%%%END
%%%BLOCK id=101-03 ch=15 sec=折射定律和全反射规律的综合应用 kind=summary alt=none cont=none y=0.32-0.40
③\ 折射定律和全反射规律的综合应用
\[ \theta\ge C=\arcsin\frac{1}{n}\qquad v=\frac{c}{n}\qquad t=\frac{L}{v} \]
光路控制、光的色散。
%%%END
%%%BLOCK id=101-04 ch=15 sec=平行玻璃砖·圆柱体·三棱镜 kind=summary alt=none cont=none y=0.42-0.45
④\ 平行玻璃砖\quad 圆柱体、三棱镜
%%%END
%%%BLOCK id=101-05 ch=15 sec=光的波动性·电磁波·相对论 kind=summary alt=14 cont=none y=0.45-0.62
⑤\ 光的波动性、电磁波、相对论
% 原稿：此标题右侧一个大括号括出下面三项
\begin{itemize}
\item 光的\uline{干涉}、衍射、和偏振。
  \begin{itemize}
  \item 加强减弱的区交替。
  \item $\uparrow$光波长略大或与障碍小孔尺寸相近
    % 原稿此项与上一项写在同一行；“↑”箭头指向上方的“衍射”；其右侧有大括号括出以下三项
    \begin{itemize}
    \item 单缝衍射
    \item 圆孔衍射
    \item 圆盘衍射。
    \end{itemize}
  \end{itemize}
\item 电磁波
\item 相对论。
\end{itemize}
%%%END
%%%BLOCK id=101-06 ch=15 sec=考点·光学与电磁波相对论 kind=summary alt=14 cont=none y=0.64-0.80
% 原稿：左缘有一个“L”形标记（疑为序号⑥，未画圈）；页面左缘 y≈0.86 处另有半个圆圈（疑为序号⑦的圈，内无内容）
\unclear{⑥}\ \textbf{考点。}
% 以下三项在原稿中被“考点”右侧的大括号括起
\begin{itemize}
\item 光的干涉现象。
\item 光的衍射和偏振现象
\item 电磁波与相对论
\end{itemize}
%%%END
%%%PAGEEND 101
